% Paper 3 — Discussion Section
% Building Local-Language Economic Narrative Indices

\section{Discussion}\label{sec:discussion}

\subsection{\texorpdfstring{The $\kappa$ Paradox}{The kappa Paradox}: Why High Accuracy Does Not Imply Reliable Classification}

On the frozen BENI v1 validation set, the repaired TF-IDF model achieves 88.17\% accuracy, 0.736 $\kappa$, and 0.823 F1 on the Economic class. The majority-class baseline still reaches 62.90\% accuracy, which is why accuracy alone remains an incomplete metric on imbalanced text classification: the more important question is how well the model recovers the minority Economic class.

The practical implication is that models must be evaluated on $\kappa$ or F1 for the minority class, not accuracy. The TF-IDF model's F1 of 0.823 on the frozen BENI v1 validation set is the operative production baseline for the repaired release. Researchers building similar indices should pre-register F1 (Economic class) or $\kappa$ as their primary metric, not accuracy.

\subsection{Why Short-Run Correlations Are Null}

The strong level correlations ($-$0.72 to $-$0.75) and null first-differenced correlations ($-$0.04 to 0.10) in the calibration experiment suggest that economic narrative salience in Bangladeshi news may operate primarily as a \textit{long-run structural} measurement signal rather than a \textit{short-run cyclical} forecasting signal.

Three interpretations are consistent with this pattern:

\begin{enumerate}
    \item \textbf{Structural transformation}: Between 2014 and 2020, Bangladesh's economy underwent significant structural change: income rose, the garment sector expanded, and digital financial services proliferated. The declining share of economic news in the calibration index may reflect a genuine shift in news composition as the economy diversified beyond headline economic stories.
    \item \textbf{News production dynamics}: Bangladeshi newspapers may allocate front-page space to economic narratives only during periods of acute economic stress or major policy events, producing a low-frequency signal that correlates with the level of macroeconomic indicators but not their monthly changes.
    \item \textbf{Nominal vs. real signal}: The level correlations may partly reflect nominal trends (both the exchange rate and consumer prices rose steadily over the period), and differencing removes this shared trend, revealing the absence of a short-run relationship.
\end{enumerate}

This null result is a methodological contribution in itself: it demonstrates that \textbf{classification improvement does not guarantee correlation improvement}. A hypothetical model that achieves very high F1 on economic relevance classification would not necessarily produce first-differenced correlations above $|r| = 0.10$, because the mapping from classification accuracy to time-series predictive power is mediated by the frequency, magnitude, and timing of narrative shifts. This finding echoes the broader sentiment forecasting literature, where publication bias inflates reported correlations~\citep{nabil2026economic}. For BENI v1, the practical implication is that the final source-balanced index must be validated separately after regeneration from the harmonised database.

\subsection{Qualitative Error Analysis}

The LLM annotator's self-consistency evaluation (Section~\ref{sec:methods}) identified 2 disagreements in 50 articles ($\kappa = 0.48$). Both cases illuminate the boundary of what constitutes ``economic relevance'' for BENI:

\begin{itemize}
    \item \textbf{Political party funding} (article \texttt{bnlp\_dev\_1208}): A debate about funding sources for political campaigns was classified as Not Economic on the first pass (focusing on the political dimension) but Economic on the second pass (focusing on the fiscal policy implications of campaign finance). The article contains economic vocabulary (budget, funding, allocation) but the primary frame is political. This suggests that borderline articles occupy a genuine grey zone where even consistent annotators disagree.
    \item \textbf{Mobile phone discounts} (article \texttt{bnlp\_train\_6027}): A report on promotional pricing by mobile operators was classified as Not Economic initially (consumer pricing, not macro-economics) but Economic on re-annotation (trade/commerce implications). The article discusses market competition and pricing strategies---microeconomic phenomena that border on trade and commerce for the BENI framework.
\end{itemize}

These disagreements are not annotator errors but genuine ambiguity at the classification boundary. They represent the irreducible uncertainty inherent in the ``Economic'' label, which no amount of annotation budget expansion or model improvement can eliminate. For the BENI pipeline, we estimate this irreducible disagreement rate at approximately 4\% (2/50), corresponding to a self-consistency ceiling of $\kappa \approx 0.48$.

\subsection{Implications for Methodology}

Three findings from this calibration study have direct implications for researchers building local-language economic narrative indices:

\begin{enumerate}
    \item \textbf{Frozen validation clarifies the production baseline.} The repaired TF-IDF model performs substantially better than the majority baseline on the frozen 186-row BENI v1 validation set, so it is the production comparator for the current release. The legacy BanglaBERT calibration result remains useful as prototype evidence, but it is not part of the frozen release tables.
    \item \textbf{Accuracy is a misleading metric for imbalanced economic text.} With base rates of 5--10\% (typical for economic relevance classification), accuracy is dominated by majority-class performance. Researchers should report $\kappa$ or minority-class F1 as their primary evaluation metric.
    \item \textbf{Active learning with linear classifiers has limited value below $k = 300$ annotations at low base rates.} The flat active learning curve shows that TF-IDF cannot learn the minority class from fewer than 300 articles with a 7.4\% positive rate. At these sample sizes, the annotation budget is better spent on acquiring higher-quality labels (e.g., through LLM annotation or uncertainty-guided sampling) than on training increasingly complex models on the same small set.
\end{enumerate}

\subsection{Implications for BENI v1}

The BENI v1 database changes the scale and purpose of the pipeline. The calibration experiment establishes the measurement logic, but the final BENI v1 index should be generated from the harmonised 1.47-million-record database with explicit source balancing, duplicate handling, and versioned model outputs. This matters because a high-volume source can dominate raw article-weighted monthly averages, and because post-2020 BNAD coverage extends the period into a different macroeconomic regime.

The next methodological step is therefore mechanical but important: freeze the annotation set, choose the production classifier, run predictions over the deduplicated BENI v1 article file, construct article-weighted and source-balanced monthly indices, and rerun the validation tables. Only after that step should the index be used in the subsequent macroeconomic validation or inflation-nowcasting papers.
